% !TeX program = lualatex
% !TeX encoding = utf8
% !TeX spellcheck = uk_UA
% !TeX root = ../ElMagConspect.tex

%=========================================================
\chapter{Електромагнітні хвилі}\label{\currfilebase}
\graphicspath{{\currfilebase/Pictures}}
%=========================================================
\epigraph{Поле не лише переносить взаємодію між зарядами і не лише зберігає енергію --- у відсутності будь-яких зарядів воно здатне жити
	цілком самостійним життям, біжучи крізь порожній простір зі швидкістю світла. Це і є, мабуть, найпереконливіший доказ того, що поле --- не
	математична абстракція, а такий самий фізичний об'єкт, як і речовина.
}{Переосмислення ідей сучасної фізики}

%% --------------------------------------------------------
\section{Хвильове рівняння}
%% --------------------------------------------------------

\subsection*{Хвильовий процес}

Попри величезну різноманітність фізичних процесів, що спричиняють хвилі (коливання струни чи ланцюжка зв'язаних тягарців, хвилі на поверхні
води, звукові хвилі, електромагнітні хвилі), їхнє утворення підпорядковане спільному принципу: збурення, що виникло в деякий момент часу в
одному місці, з часом поширюється в просторі з певною швидкістю. \emphx{Хвиля --- це процес поширення збурень у просторі}.

Хвилю описують деякою величиною $u = f(x,y,z,t)$, що є функцією як просторових координат, так і часу. Наприклад, для хвилі на поверхні води $u$
--- це висота підняття поверхні над рівноважним рівнем; для механічної хвилі в ланцюжку зв'язаних пружинками тягарців (рис.~\ref{pic:Chain}) $u$
--- зміщення тягарця від положення рівноваги; для електромагнітної хвилі $u$ --- змінна складова електричного чи магнітного поля, що
накладається на їхнє (можливо, нульове) статичне значення.

\begin{figure}[h!]\centering
	\localinput{chain.tikz}
	\caption{Механічна хвиля в ланцюжку зв'язаних пружинками тягарців --- наочна модель для
		розуміння хвильового процесу\label{pic:Chain}}
\end{figure}

\subsection*{Хвильове рівняння}

Розглянемо збурення, що поширюється вздовж осі $x$ без зміни форми зі швидкістю $v$ (рис.~\ref{pic:Perturbation}). Якщо в момент $t=0$ профіль
збурення описувався функцією $u=f(x)$, то через час $t$ те саме за формою збурення просто зміститься на відстань $vt$, тобто опишеться функцією
$u = f(x-vt)$ (для хвилі, що біжить у додатному напрямку осі $x$) або $u=g(x+vt)$ (для хвилі, що біжить у від'ємному напрямку). У загальному
випадку (суперпозиція обох):
\begin{equation}\label{eq:DAlembertSolution}
	u = f(x-vt) + g(x+vt).
\end{equation}

\begin{figure}[h!]\centering
	\localinput{perturbation.tikz}
	\caption{Збурення, що поширюється без зміни форми зі швидкістю $v$: профіль у момент $t$ ---
		це профіль у момент $0$, зсунутий на $vt$\label{pic:Perturbation}}
\end{figure}

Конкретний вигляд функцій $f$ і $g$ визначається початковими умовами (початковою формою збурення й початковим розподілом швидкостей) і може бути
яким завгодно. Знайдемо диференціальне рівняння, якому підпорядковується \emphx{будь-яка} функція вигляду~\eqref{eq:DAlembertSolution},
незалежно від її конкретного вигляду. Позначимо $\xi = x-vt$; тоді, за правилом диференціювання складеної функції,
\begin{equation*}
	\frac{\partial u}{\partial x} = f'(\xi), \qquad \frac{\partial^2 u}{\partial x^2} = f''(\xi); \qquad \frac{\partial u}{\partial t} = -vf'(\xi),
	\qquad \frac{\partial^2u}{\partial t^2} = v^2f''(\xi)
\end{equation*}
(і аналогічно для доданка $g(x+vt)$, з тим самим кінцевим результатом через парність других похідних за знаком швидкості). Порівнюючи другі
похідні, знаходимо \emphx{хвильове рівняння}:
\begin{equation}\label{eq:WaveEquation1D}
	\tcbhighmath{
		\frac{\partial^2 u}{\partial x^2} = \frac1{v^2}\frac{\partial^2u}{\partial t^2}.
	}
\end{equation}
У тривимірному випадку, коли збурення поширюється в довільних напрямках, хвильове рівняння узагальнюють, замінюючи другу похідну за $x$ на
лапласіан $\nabla^2 = \dfrac{\partial^2}{\partial x^2}+\dfrac{\partial^2}{\partial y^2}+\dfrac{\partial^2}{\partial z^2}$:
\begin{equation}\label{eq:WaveEquation3D}
	\tcbhighmath{
		\nabla^2 u = \frac1{v^2}\frac{\partial^2u}{\partial t^2}.
	}
\end{equation}

\begin{Attention}
	Рівняння~\eqref{eq:WaveEquation3D} --- своєрідний \enquote{пароль}, за яким фізики впізнають хвильові процеси в геть різних розділах фізики:
	\emphx{якщо якісь закони фізики (механіки, акустики, електродинаміки, квантової механіки) приводять до рівняння саме такого вигляду, то ці
		закони описують поширення хвиль} зі швидкістю $v$, яка входить у це рівняння, --- незалежно від конкретної фізичної природи величини $u$.
\end{Attention}

\subsection*{Плоска монохроматична хвиля}

\emphx{Плоскою хвилею} називають хвилю, що має однакове значення $u$ в усіх точках площини, перпендикулярної до напряму її поширення
(\emphx{хвильової площини}, рис.~\ref{pic:WavevectorGeometry}). Хвилю називають \emphx{монохроматичною}, якщо $u$ змінюється з часом за
гармонічним законом з певною частотою. Для плоскої монохроматичної хвилі:
\begin{equation}\label{eq:PlaneWaveDef}
	\tcbhighmath{
		u = u_m\cos(\omega t - \vect k\cdot\vect r).
	}
\end{equation}

\begin{figure}[h!]
	\centering
	\begin{minipage}[t]{0.48\linewidth}\centering
		\localinput{wavevector-geometry.tikz}
	\end{minipage}\hfill
	\begin{minipage}[t]{0.48\linewidth}\centering
		\includegraphics[width=\linewidth]{plane_wave}
	\end{minipage}
	\caption{Хвильовий вектор $\vect k$ вказує напрям поширення плоскої хвилі й перпендикулярний
		до хвильових площин\label{pic:WavevectorGeometry}}
\end{figure}

Для зручності плоску монохроматичну хвилю часто подають у комплексній формі:
\begin{equation}\label{eq:PlaneWaveComplex}
	u = u_m\, e^{i(\omega t - \vect k\cdot\vect r)},
\end{equation}
розуміючи, що фізичний зміст має лише \emphx{дійсна} частина цього виразу (сама ж комплексна форма --- лише зручний обчислювальний прийом:
формула Ейлера $e^{i\theta}=\cos\theta+i\sin\theta$ дає змогу компактно поєднати синус і косинус, а диференціювання й інтегрування зводяться до
простого множення на $i\omega$ чи $-i\vect k$).

\subsection*{Характеристики монохроматичної хвилі}

\emphx{Довжиною хвилі} $\lambda$ називають відстань, на яку зсунуті одна від одної дві точки простору, що коливаються цілком синхронно (в одній
фазі):
\begin{equation}\label{eq:Wavelength}
	\lambda = \frac{2\pi}{k},
\end{equation}
де $\vect k$ --- \emphx{хвильовий вектор}, напрямлений уздовж поширення хвилі перпендикулярно до хвильової площини, а $k=|\vect k|$ ---
\emphx{хвильове число}. Величину $\omega$ у формулі~\eqref{eq:PlaneWaveDef} називають \emphx{коловою частотою} хвилі, а $T=2\pi/\omega$ ---
\emphx{періодом}. Нарешті, \emphx{фазовою швидкістю} хвилі називають швидкість, з якою в просторі переміщується точка з фіксованою фазою
$\omega t - \vect k\cdot\vect r = \const$:
\begin{equation}\label{eq:PhaseVelocity}
	\tcbhighmath{
		v = \frac\omega k, \qquad \vect v = \frac\omega k\,\frac{\vect k}k.
	}
\end{equation}

%% --------------------------------------------------------
\section{Електромагнітні хвилі}
%% --------------------------------------------------------

\subsection*{Виведення хвильових рівнянь з рівнянь Максвелла}

Розглянемо електромагнітне поле в однорідній ділянці простору, де немає вільних зарядів і струмів провідності ($\rho=0$, $\vect j=0$). Рівняння
Максвелла (гл.~\ref{MaxwellEqns}) у цьому випадку спрощуються:
\begin{equation}\label{eq:MaxwellNoSources}
	\divg\Bfield=0, \qquad \rot\Efield = -\frac1c\parttime{\Bfield}, \qquad \divg\Dfield=0, \qquad \rot\Hfield = \frac1c\parttime{\Dfield},
\end{equation}
доповнені матеріальними рівняннями $\Dfield=\epsilon\Efield$, $\Bfield=\mu\Hfield$ (нехай $\epsilon,\mu=\const$).

Застосуємо операцію $\rot$ до другого рівняння:
\begin{equation*}
	\rot\rot\Efield = -\frac1c\parttime{}\left(\rot\Bfield\right) = -\frac\mu{c}\parttime{}\left(\rot\Hfield\right) = -\frac{\epsilon\mu}{c^2}\pparttime{\Efield}
\end{equation*}
(в останній рівності скористались четвертим рівнянням). З іншого боку, за тотожністю векторного аналізу й $\divg\Efield=0$ (яке випливає з
$\divg\Dfield=0$): $\rot\rot\Efield = \grad(\divg\Efield)-\nabla^2\Efield = -\nabla^2\Efield$. Прирівнюючи, отримуємо хвильове рівняння для
електричного поля:
\begin{equation}\label{eq:WaveEqE}
\tcbhighmath{
	\nabla^2\Efield = \frac1{v^2}\pparttime{\Efield}.
}
\end{equation}
Цілком аналогічно, застосувавши $\rot$ до четвертого рівняння й підставивши друге, отримуємо те саме рівняння для магнітного поля:
\begin{equation}\label{eq:WaveEqH}
	\tcbhighmath{
		\nabla^2\Hfield = \frac1{v^2}\pparttime{\Hfield}, \qquad v = \frac{c}{\sqrt{\epsilon\mu}}.
	}
\end{equation}

\subsection*{Швидкість поширення і показник заломлення}

Отже, поля $\Efield$ і $\Hfield$ підпорядковуються \emphx{хвильовому рівнянню}, тобто описують процес поширення збурень електромагнітного поля
--- \emphx{електромагнітну хвилю}, --- зі швидкістю
\begin{equation}\label{eq:WaveSpeed}
	v = \frac{c}{\sqrt{\epsilon\mu}}.
\end{equation}
У вакуумі ($\epsilon=\mu=1$) швидкість поширення дорівнює електродинамічній сталій $c\approx3\cdot10^{10}$~см/с --- \emphx{швидкості світла}. У
середовищі з $\epsilon,\mu>1$ швидкість поширення менша, $v<c$, на величину
\begin{equation}\label{eq:RefractiveIndex}
	\tcbhighmath{
		n = \sqrt{\epsilon\mu},
	}
\end{equation}
яку називають \emphx{абсолютним показником заломлення} середовища.

\begin{Attention}
	Той факт, що рівняння Максвелла --- сформульовані без жодного посилання на оптику чи хвилі, виключно на основі дослідів з електрикою й
	магнетизмом (закон Кулона, закон Біо--Савара--Лапласа, закон індукції, струм зміщення) --- \emphx{самі собою} приводять до хвильового
	рівняння зі швидкістю, що чисельно збігається зі швидкістю світла, --- один із найграндіозніших результатів усієї класичної фізики. Саме на
	цій підставі Максвелл висунув гіпотезу, що \emphx{світло --- це і є електромагнітна хвиля}, а оптика --- лише розділ електродинаміки.
\end{Attention}

\subsection*{Плоскі електромагнітні хвилі}

Розглянемо гармонічні електромагнітні хвилі в комплексній формі:
\begin{equation}\label{eq:PlaneEMWave}
	\vect E = \vect E_m\, e^{i(\omega t-\vect k\cdot\vect r)}, \qquad \vect H = \vect H_m\, e^{i(\omega t-\vect k\cdot\vect r)}.
\end{equation}
Підставляючи ці вирази в рівняння Максвелла, диференціальні операції заміняються алгебраїчними за правилом
\begin{equation}\label{eq:OperatorSubstitution}
	\grad \to -i\vect k, \qquad \parttime{} \to i\omega
\end{equation}
(легко переконатися прямим диференціюванням~\eqref{eq:PlaneEMWave}). Рівняння Максвелла~\eqref{eq:MaxwellNoSources} набувають вигляду
\begin{equation}\label{eq:MaxwellPlaneWave}
	\vect k\cdot\Efield=0, \qquad \vect k\cdot\Bfield=0, \qquad \vect k\times\Efield = \frac\omega{c}\Bfield, \qquad -\vect k\times\Hfield =
	\frac\omega c\epsilon\Efield.
\end{equation}

\subsection*{Поперечність електромагнітних хвиль}

Перші два рівняння в~\eqref{eq:MaxwellPlaneWave} показують, що
\begin{equation}\label{eq:Transversality}
	\tcbhighmath{
		\Efield \perp \vect k, \qquad \Hfield \perp \vect k:
	}
\end{equation}
вектори $\Efield$ і $\Hfield$ перпендикулярні до напряму поширення хвилі. Цю властивість називають \emphx{поперечністю} електромагнітних хвиль.
Два інші рівняння дають співвідношення для амплітуд (і показують, що $\Efield\perp\Hfield$, тобто $\Efield$, $\Hfield$, $\vect k$ утворюють
правовинтову трійку взаємно перпендикулярних векторів):
\begin{equation}\label{eq:AmplitudeRelation}
	\tcbhighmath{
		nE_m = B_m, \qquad \sqrt\epsilon\, E_m = \sqrt\mu\, H_m.
	}
\end{equation}

\begin{Attention}
	Поперечність --- принципова властивість, що відрізняє електромагнітні хвилі, наприклад, від звукових (поздовжніх). Разом
	з~\eqref{eq:AmplitudeRelation} вона означає, що в біжучій електромагнітній хвилі $\Efield$ і $\Hfield$ не є незалежними: досить знати одне з
	них (наприклад, $\Efield$) у кожен момент часу, щоб миттєво відновити й друге.
\end{Attention}

%% --------------------------------------------------------
\section{Енергія та вектор Пойнтінга для хвилі}
%% --------------------------------------------------------

Густина енергії електромагнітного поля (гл.~\ref{MaxwellEqns}) для амплітудних значень:
\begin{equation}\label{eq:EnergyDensityAmplitude}
	w_m = \frac{\epsilon E_m^2}{8\pi}+\frac{\mu H_m^2}{8\pi}.
\end{equation}
Скориставшись співвідношенням амплітуд~\eqref{eq:AmplitudeRelation}, $\sqrt{\epsilon/\mu}\,E_m = H_m$, електрична й магнітна частини густини
енергії виявляються \emphx{рівними} одна одній, і
\begin{equation*}
	w_m = \frac{\epsilon E_m^2}{4\pi} = \frac{\mu H_m^2}{4\pi}.
\end{equation*}

Знайдемо вектор Пойнтінга $\vect\Pi = \dfrac{c}{4\pi}\left[\Efield\times\Hfield\right]$ для хвилі. Скориставшись зв'язком
$\Efield = -\dfrac{c}{\epsilon\omega}\left[\vect k\times\Hfield\right]$ (з останнього рівняння~\eqref{eq:MaxwellPlaneWave}) і тотожністю подвійного
векторного добутку:
\begin{equation}\label{eq:PoyntingWaveDerivation}
	\vect\Pi = \frac{c^2}{4\pi\epsilon\omega}\left[\left[\Hfield\times\vect k\right]\times\Hfield\right] = \frac{c^2}{4\pi\epsilon\omega}\,H^2\vect
	k = \frac{c}{\omega}\cdot\frac1{\epsilon\mu}\left(\frac{\mu H^2}{4\pi}\right)c\vect k.
\end{equation}
Помічаючи, що $\mu H^2/(4\pi) = w$ (густина енергії, з точністю до множника $2$ від переходу амплітуда$\to$миттєве значення, який скорочується
однаково в обох частинах), а $c/(\epsilon\mu) \cdot c\vect k/\omega = (c/\sqrt{\epsilon\mu})(\vect k/k) = \vect v$ (за
означенням~\eqref{eq:PhaseVelocity} і~\eqref{eq:WaveSpeed}), отримуємо надзвичайно простий і красивий результат:
\begin{equation}\label{eq:PoyntingWaveResult}
	\tcbhighmath{
		\vect\Pi = w\vect v.
	}
\end{equation}

\begin{Attention}
	Формула~\eqref{eq:PoyntingWaveResult} має простий фізичний зміст: енергія електромагнітної хвилі \emphx{переноситься разом з нею} зі
	швидкістю поширення самої хвилі $\vect v$ --- вектор Пойнтінга просто дорівнює густині енергії, помноженій на швидкість її \enquote{течії}, за
	повною аналогією з потоком речовини (густина $\times$ швидкість).
\end{Attention}

Оскільки реальні (фізичні) поля --- високочастотні, для практики зазвичай цікавий не миттєвий, а \emphx{усереднений} за період коливань потік
енергії. Усереднене значення модуля вектора Пойнтінга називають \emphx{інтенсивністю} випромінювання:
\begin{equation}\label{eq:Intensity}
	I = \left\langle|\vect\Pi|\right\rangle \propto E_m^2.
\end{equation}

\begin{figure}[h!]\centering
	\localinput{emwave.tikz}
	\caption{Плоска монохроматична електромагнітна хвиля: взаємно перпендикулярні коливання
		$\Efield$ і $\Hfield$, що поширюються вздовж $\vect k$\label{pic:EMWave}}
\end{figure}

%% --------------------------------------------------------
\subsection*{Тиск електромагнітної хвилі}
%% --------------------------------------------------------

Електромагнітні хвилі, відбиваючись від тіл чи поглинаючись ними, чинять на них тиск. Цей тиск виникає внаслідок дії магнітного поля хвилі на
електричні струми, збуджені в тілі електричним полем тієї самої хвилі (рис.~\ref{pic:EMPressure}).

\begin{figure}[h!]\centering
	\localinput{EMpresure.tikz}
	\caption{Механізм виникнення тиску електромагнітної хвилі: сила Ампера $d\vect f$, що діє на
		струм, збуджений полем $\Efield$, у полі $\Bfield$\label{pic:EMPressure}}
\end{figure}

Електричне поле хвилі збуджує в тілі струм густиною $\vect j=\lambda\Efield$, на який з боку магнітного поля хвилі $B=nE=(c/v)E$ діє сила Ампера
$d\vect F = \tfrac1c[\vect j\times\Bfield]\,dV$, спрямована вздовж поширення хвилі. Якщо вся енергія, що падає на поверхню тіла ($\Pi\, dS = wv\,
dS$), перетворюється на теплоту ($\vect j\cdot\Efield\, dV$), тобто $wv\,dS = jE\,dV$, то сила, що діє на поверхню:
\begin{equation}\label{eq:RadiationForce}
	dF = \frac1c jB\,dV = \frac1c j\frac{c}{v}E\,dV = \frac1v jE\,dV = \frac1v wv\,dS = w\,dS,
\end{equation}
звідки тиск при повному поглинанні
\begin{equation}\label{eq:RadiationPressureFull}
	\tcbhighmath{
		p = w.
	}
\end{equation}
У випадку часткового відбивання (коефіцієнт відбивання $0\leq\rho\leq1$) тіло отримує додатковий імпульс від відбитої хвилі, і тиск
\begin{equation}\label{eq:RadiationPressurePartial}
	\tcbhighmath{
		p = w(1+\rho).
	}
\end{equation}

\begin{Attention}
	При повному відбиванні ($\rho=1$) тиск рівно \emphx{вдвічі} більший, ніж при повному поглинанні ($\rho=0$), --- саме цю різницю в $1899$~р.
	і виміряв П.~М.~Лебедєв у своєму класичному досліді (рис.~\ref{pic:LebedevExp}): крутильні терези з двома крильцями --- зачорненим
	(поглинальним) і дзеркальним (відбивальним) --- поверталися моментом сил, зумовленим саме тим, що дзеркальне крильце отримувало вдвічі
	більший тиск світла, ніж зачорнене.
\end{Attention}

\begin{figure}[h!]\centering
	\includegraphics[width=0.65\linewidth]{LebedevExp}
	\caption{Схема досліду П.~М.~Лебедєва (1899~р.) з вимірювання тиску світла\label{pic:LebedevExp}}
\end{figure}

\subsection*{Тиск світла в природі та техніці}

Тиск сонячного світла, попри свою мізерну величину (на орбіті Землі --- близько \(5\cdot10^{-6}\)~Па),
відіграє помітну роль і в природних явищах, і в сучасних космічних технологіях.

В астрофізиці тиск випромінювання, поряд із тиском газу, протидіє гравітаційному стисканню
й забезпечує стійкість зір. У масивних гарячих зорях він стає настільки суттєвим, що
визначає їхню структуру та еволюцію. Саме тиском сонячного світла (разом із сонячним вітром)
пояснюється форма кометних хвостів: дрібні пилинки й молекули газу «здуваються» в напрямку
\emphx{від} Сонця, тому хвіст комети завжди орієнтований приблизно протилежно до Сонця,
незалежно від напряму руху самої комети.

Ідея \emphx{сонячного вітрила} полягає в тому, щоб використати тиск світла для розгону
космічного апарата без витрати палива. Першу спробу випробувати вітрило в космосі здійснила
Планетарна спілка у 2005~році (місія \emphx{Cosmos~1}), але ракета-носій зазнала аварії,
і апарат не досяг орбіти. Лише у 2019~році місія \emphx{LightSail~2} тієї ж організації
вперше продемонструвала кероване сонячне вітрилоплавання: апарат успішно підняв свою орбіту,
використовуючи виключно тиск сонячних фотонів.

Значно амбітнішим є проєкт \emphx{Breakthrough Starshot}, анонсований у 2016~році.
Його мета --- розігнати грамові зонди на лазерному вітрилі до швидкості \(0{,}2c\)
за допомогою наземної системи лазерів потужністю близько \(100\)~ГВт, щоб досягти
системи Альфи Центавра приблизно за \(20\)~років. Важливо, що це \emphx{пролітна} місія:
зонд не гальмує і не виходить на орбіту навколо мети, а лише короткочасно спостерігає
систему на швидкості \(0{,}2c\). Станом на 2026~рік проєкт не запустив жодного зонда;
його фактичне фінансування обмежилося кількома мільйонами доларів у вигляді мікрогрантів.

Отже, практичного застосування тиску світла в промисловому сенсі поки що немає ---
це радше \emphx{перспективна технологія} та \emphx{природний чинник}, важливий для
розуміння будови зір і кометних хвостів.

%\begin{figure}[h!]
%	\centering
%	\begin{minipage}[t]{0.45\linewidth}\centering
%		\includegraphics[width=\linewidth]{Cosmos_solar_sail}\\
%		\footnotesize Сонячне вітрило (проєкт \enquote{Космос})
%	\end{minipage}\hfill
%	\begin{minipage}[t]{0.45\linewidth}\centering
%		\includegraphics[width=\linewidth]{Starshot-light-sail_Breakthroughinitiatives}\\
%		\footnotesize Концепція проєкту Breakthrough Starshot
%	\end{minipage}
%	\caption{Сонячні та лазерні вітрила\label{pic:SolarSail}}
%\end{figure}


%\begin{figure}[h!]
%	\centering
%	\begin{minipage}[t]{0.45\linewidth}\centering
%		\includegraphics[width=\linewidth]{star}
%	\end{minipage}\hfill
%	\begin{minipage}[t]{0.45\linewidth}\centering
%		\includegraphics[width=\linewidth]{Comet}
%	\end{minipage}
%	\caption{Тиск світла в астрофізиці: рівновага зорі (ліворуч) і хвіст комети (праворуч),
%		спрямований від Сонця\label{pic:StarComet}}
%\end{figure}

%% --------------------------------------------------------
\section{Імпульс і момент імпульсу електромагнітного поля}
%% --------------------------------------------------------

\subsection*{Густина імпульсу поля}

Оскільки поле має енергію, а електромагнітна хвиля чинить тиск на речовину (передаючи їй імпульс), логічно припустити, що й саме поле, в одиниці
об'єму, має певний \emphx{імпульс}. Для електромагнітної хвилі у вакуумі $\Pi = wc$, тому $\Pi/c^2 = w/c$. З релятивістського співвідношення
енергії й імпульсу, $E^2 = m^2c^4+p^2c^2$, для частинки з нульовою масою спокою (наприклад, фотона) $E=pc$, тобто $p=E/c$. Порівнюючи
структуру цих двох формул, для густини імпульсу поля природно отримуємо
\begin{equation}\label{eq:MomentumDensityDef}
	\tcbhighmath{
		\vect g = \frac{\vect\Pi}{c^2}.
	}
\end{equation}

Переконаємось у цьому прямим розрахунком, повторюючи міркування про тиск хвилі. Сила Ампера $d\vect f = \tfrac1c[\vect j\times\Bfield]$, що діє
на струм, збуджений хвилею в тілі, за час $dt$ змінює імпульс одиниці об'єму тіла на величину
\begin{equation*}
	d\vect g = \vect f\, dt = \frac1c\left[\vect j\times\Bfield\right]dt = \frac\lambda c\left[\Efield\times\Bfield\right]dt =
	\frac{4\pi\lambda}{c^2}\vect\Pi\, dt = \frac{4\pi\lambda}c\, w\vect n\, dt,
\end{equation*}
тоді як за той самий час у тій самій одиниці об'єму поглинається (передається від поля речовині) енергія
\begin{equation*}
	dw = \vect j\cdot\Efield\, dt = \lambda E^2\, dt = 4\pi\lambda\frac{E^2}{4\pi}\, dt = 4\pi\lambda w\, dt.
\end{equation*}
Звідси $d\vect g = \dfrac{dw}{c}\vect n$, тобто за весь час дії хвилі
\begin{equation}\label{eq:MomentumDerivation}
	\vect g = \frac{w}c\vect n = \frac{\vect\Pi}{c^2}
\end{equation}
--- прямий розрахунок підтверджує результат~\eqref{eq:MomentumDensityDef}, отриманий раніше з релятивістської аналогії.

\begin{Attention}
	Погодженість двох зовсім різних способів отримання формули~\eqref{eq:MomentumDensityDef} --- прямого мікроскопічного розрахунку сили Ампера
	й релятивістського співвідношення енергія--імпульс для безмасової частинки --- ще один вияв внутрішньої узгодженості класичної
	електродинаміки з рештою фізики (тут --- зі спеціальною теорією відносності, яка історично якраз і \emphx{виникла} з аналізу рівнянь
	Максвелла).
\end{Attention}

\subsection*{Приклад: момент імпульсу поля циліндричного конденсатора (парадокс Фейнмана)}

Заряджений циліндричний конденсатор (зовнішня обкладка радіуса $R$, внутрішньою знехтуємо) вміщений у зовнішнє однорідне магнітне поле
$\Bfield$, напрямлене вздовж його осі; конденсатор може вільно обертатись навколо цієї осі. Заряд конденсатора $q$, маса --- $m$
(рис.~\ref{pic:CapacitorAngularMomentum}). Знайдемо момент імпульсу поля конденсатора та кутову швидкість, з якою конденсатор почне обертатися
після вимкнення поля $\Bfield$.

\begin{figure}[h!]\centering
	\localinput{capacitor-angular-momentum.tikz}
	\caption{Циліндричний конденсатор в аксіальному магнітному полі: поля $\Efield$, $\Bfield$ і
		вектор Пойнтінга $\vect\Pi$, що циркулює навколо осі\label{pic:CapacitorAngularMomentum}}
\end{figure}

Густина моменту імпульсу поля відносно осі:
\begin{equation}\label{eq:AngularMomentumDensity}
	\vec{\mathcal L} = \vect r\times\frac{\vect\Pi}{c^2} = \frac1{4\pi c}\vect r\times\left[\Efield\times\Bfield\right] = -\frac1{4\pi
	c}\Bfield\left(\vect r\cdot\Efield\right)
\end{equation}
(тут скористались тотожністю подвійного векторного добутку й тим, що $\Bfield\perp\vect r$). Поле між обкладками конденсатора, за теоремою
Гаусса, $E = 2q/(rh)$, де $h$ --- висота циліндра. Інтегруючи модуль~\eqref{eq:AngularMomentumDensity} за об'ємом між обкладками:
\begin{equation}\label{eq:CapacitorAngularMomentum}
	\tcbhighmath{
		L = \iiint\limits_V \mathcal L\, dV = \frac1{4\pi c}B\,\frac{2q}h\,\pi R^2 h = \frac{qBR^2}{2c}, \qquad \vect L = -\frac{qR^2}{2c}\Bfield.
	}
\end{equation}

При вимкненні магнітного поля момент імпульсу, який раніше \enquote{зберігався} в самому електромагнітному полі, нікуди не зникає --- за законом
збереження моменту імпульсу він переходить у механічне обертання самого конденсатора, $L = mR^2\omega$ (наближено вважаючи всю масу
розподіленою на радіусі $R$), звідки
\begin{equation}\label{eq:CapacitorOmega}
	\tcbhighmath{
		\vect\omega = -\frac{q}{2mc}\Bfield.
	}
\end{equation}

\begin{Attention}
	Ця задача --- варіант відомого \emphx{парадоксу Фейнмана} (\emph{Feynman's disk paradox}): до вимкнення поля конденсатор \emphx{нерухомий}, і
	жодного механічного моменту імпульсу у видимій системі немає --- проте \emphx{момент імпульсу самого електромагнітного поля}, обчислений за
	формулою~\eqref{eq:CapacitorAngularMomentum}, відмінний від нуля! Коли поле $\Bfield$ вимикають, змінне магнітне поле (гл.~\ref{ElMagInduction})
	породжує вихрове електричне поле, яке й розкручує заряджені обкладки конденсатора, --- момент імпульсу \enquote{перетікає} з поля в
	механічне обертання речовини, а повний момент імпульсу системи \emphx{поле + речовина} зберігається протягом усього процесу. Це переконливий
	приклад того, що електромагнітне поле --- цілком повноцінний фізичний об'єкт, здатний нести не лише енергію, а й імпульс, і момент імпульсу,
	--- так само, як і звичайна речовина.
\end{Attention}

%% --------------------------------------------------------
\section{Відбивання хвиль від поверхні двох діелектриків}
%% --------------------------------------------------------

\subsection*{Поляризація хвиль}

\emphx{Поляризація} --- характеристика поперечних хвиль, що описує поведінку вектора величини, яка коливається (для електромагнітної хвилі ---
вектора $\Efield$), у площині, перпендикулярній до напряму поширення. При розгляді відбивання й заломлення хвилі на межі розділу двох середовищ
зручно окремо розглядати два незалежні випадки поляризації відносно \emphx{площини падіння} (площини, що містить хвильовий вектор падаючої хвилі
й нормаль до межі розділу):
\begin{itemize}
	\item \emphx{$s$-поляризація} (від нім. \emph{senkrecht} --- перпендикулярний): вектор $\Efield$ перпендикулярний до площини падіння (а
	      отже, паралельний до межі розділу середовищ);
	\item \emphx{$p$-поляризація} (від \emph{parallel} --- паралельний): вектор $\Efield$ лежить у площині падіння.
\end{itemize}
Довільну поляризацію хвилі завжди можна розкласти на суперпозицію цих двох незалежних випадків.

\subsection*{Відбивання й заломлення $s$-поляризованої хвилі: закони відбивання й заломлення}

Розглянемо межу розділу ($x=0$) двох діелектриків з показниками заломлення $n_1$ (при $x<0$) і $n_2$ (при $x>0$), на яку під кутом $i$ падає
плоска $s$-поляризована хвиля (рис.~\ref{pic:SPolarization}). Породжуються \emphx{відбита} хвиля (кут $i'$) і \emphx{заломлена} хвиля (кут $r$):
\begin{equation*}
	E = E_m e^{i(\omega t - k_xx-k_yy)}, \qquad E' = E'_m e^{i(\omega' t-k'_xx-k'_yy)}, \qquad E'' = E''_m e^{i(\omega''t-k''_xx-k''_yy)}.
\end{equation*}

\begin{figure}[h!]\centering
	\localinput{s-polarization-geometry.tikz}
	\caption{Падаюча, відбита й заломлена $s$-поляризовані хвилі на межі розділу двох
		діелектриків\label{pic:SPolarization}}
\end{figure}

Гранична умова для тангенціальної складової $\Efield$ ($E_{\tau_1}=E_{\tau_2}$, гл.~\ref{MaxwellEqns}) при $x=0$ мусить виконуватись у
\emphx{кожен} момент часу $t$ і в \emphx{кожній} точці $y$ межі розділу одночасно:
\begin{equation*}
	E_m\, e^{i(\omega t-k_yy)} + E'_m\, e^{i(\omega' t-k'_yy)} = E''_m\, e^{i(\omega'' t-k''_yy)}.
\end{equation*}
Диференціюючи цю тотожність за $t$, а потім за $y$ (або, що те саме, порівнюючи показники експонент, які мусять збігатися при будь-яких $t,y$),
отримуємо
\begin{equation}\label{eq:BoundaryConditionE}
	\omega = \omega' = \omega'', \qquad k_y = k'_y = k''_y,
\end{equation}
після чого гранична умова спрощується до $E_m + E'_m = E''_m$. Рівність частот означає, що колір (частота) світла не змінюється при відбиванні
чи заломленні. З рівності $k_y = k'_y$ (з огляду на те, що падаюча й відбита хвилі перебувають в одному середовищі, $k=k'=n_1\omega/c$, а
$k_y = k\sin i$, $k'_y=k'\sin i'$):
\begin{equation}\label{eq:SnellReflection}
	\tcbhighmath{
		i = i'
	}
\end{equation}
--- \emphx{закон відбивання}, добре відомий з геометричної оптики. З рівності $k'_y=k''_y$ (тепер $k'=n_1\omega/c$, $k''=n_2\omega/c$):
\begin{equation}\label{eq:SnellRefraction}
	\tcbhighmath{
		n_1\sin i = n_2\sin r, \qquad \frac{\sin i}{\sin r} = \frac{n_2}{n_1} = n_{21}
	}
\end{equation}
--- \emphx{закон заломлення} (закон Снелла), де $n_{21}$ --- відносний показник заломлення другого середовища відносно першого. Оскільки
більшість діелектриків не проявляють магнітних властивостей ($\mu_1=\mu_2\approx1$),
\begin{equation}\label{eq:RelativeIndex}
	n_{21} = \sqrt{\frac{\epsilon_2}{\epsilon_1}} = \frac{v_1}{v_2}.
\end{equation}

\begin{Attention}
	Класичні закони геометричної оптики --- закони відбивання й заломлення, відомі задовго до Максвелла й зазвичай постульовані окремо (чи
	виведені з принципу Ферма) --- виявляються \emphx{прямим наслідком} граничних умов для електромагнітного поля на межі розділу середовищ.
	Максвеллівська електродинаміка не просто узгоджується з геометричною оптикою --- вона \emphx{містить} її як частинний, найпростіший випадок.
\end{Attention}

\subsection*{Формули Френеля для $s$-поляризації}

Крім уже отриманої граничної умови для $\Efield$ (яка звелась до $E_m+E'_m=E''_m$), потрібна ще одна --- гранична умова для $\Hfield$. Для
$s$-поляризованої хвилі $\Efield=E\vect e_z$ напрямлений уздовж осі $z$ (перпендикулярно до площини падіння $xy$); амплітуда й напрям
магнітного поля плоскої хвилі в середовищі з показником заломлення $n$ зв'язані з електричним співвідношенням $\Bfield=n(\hat{\vect
k}\times\Efield)$, де $\hat{\vect k}=\vect k/k$ --- одиничний вектор напрямку поширення хвилі (сам хвильовий вектор $\vect k=k\hat{\vect k}$ має
модуль $k=n\omega/c$), звідки, зокрема, $B_m=nE_m$. Оскільки одиничні вектори напрямку поширення падаючої, відбитої й заломленої хвиль дорівнюють $\hat{\vect
k}=(\cos i,\sin i,0)$, $\hat{\vect k}'=(-\cos i,\sin i,0)$, $\hat{\vect k}''=(\cos r,\sin r,0)$ (вісь $x$ --- нормаль до межі, $y$ --- уздовж
межі в площині падіння), безпосереднє обчислення $\hat{\vect k}\times\vect e_z$ дає для тангенціальної (уздовж $y$) складової $\Bfield=\Hfield$
(з огляду на $\mu_1=\mu_2\approx1$):
\begin{equation*}
	H_y = -n_1E_m\cos i, \qquad H'_y = n_1E'_m\cos i, \qquad H''_y = -n_2E''_m\cos r.
\end{equation*}
Гранична умова неперервності $H_{\tau_1}=H_{\tau_2}$, тобто $H_y+H'_y=H''_y$, набуває вигляду
\begin{equation}\label{eq:FresnelSBoundaryH}
	n_1\cos i\,(E_m-E'_m) = n_2\cos r\,E''_m.
\end{equation}
Виражаючи $n_1=n_2\sin r/\sin i$ за законом Снелла~\eqref{eq:SnellRefraction} і підставляючи $E''_m=E_m+E'_m$, отримуємо
\begin{equation*}
	\sin r\cos i\,(E_m-E'_m) = \sin i\cos r\,(E_m+E'_m) \quad\Longrightarrow\quad E_m\sin(r-i) = E'_m\sin(i+r),
\end{equation*}
де використано формулу різниці синусів. Підставляючи знайдене $E'_m$ назад у $E''_m=E_m+E'_m$ й скориставшись
$\sin(i+r)-\sin(i-r)=2\cos i\sin r$, остаточно дістаємо \emphx{формули Френеля} для $s$-поляризації:
\begin{equation}\label{eq:FresnelSReflection}
	\tcbhighmath{
		E'_\perp = -\frac{\sin(i-r)}{\sin(i+r)}\, E_\perp,
	}
\end{equation}
\begin{equation}\label{eq:FresnelSTransmission}
	\tcbhighmath{
		E''_\perp = \frac{2\sin r\cos i}{\sin(i+r)}\, E_\perp.
	}
\end{equation}

\subsection*{Формули Френеля для $p$-поляризації}

Для $p$-поляризованої хвилі (рис.~\ref{pic:PPolarization}), навпаки, у площині падіння $xy$ лежить вектор $\Efield$, а $\Bfield=B\vect e_z$
напрямлений уздовж осі $z$. Обираючи напрям $\Efield$ так, щоб трійка $\Efield,\Bfield,\hat{\vect k}$ лишалась правою (тобто за тим самим
правилом $\Bfield=n(\hat{\vect k}\times\Efield)$, як і для $s$-поляризації), одиничний вектор $\Efield$ дорівнює $\hat{\vect e}_E=\vect
z\times\hat{\vect k}$, що для падаючої, відбитої й заломленої хвиль дає тангенціальні (уздовж $y$) складові
\begin{multline*}
	E_y = E_m\cos i,\quad E'_y=-E'_m\cos i,\quad E''_y=E''_m\cos r; \\ B_z=n_1E_m,\quad B'_z=n_1E'_m,\quad B''_z=n_2E''_m
\end{multline*}
(складова $B_z$ тут цілком тангенціальна, оскільки $\Bfield\parallel\vect e_z$). Неперервність тангенціальних складових $\Efield$ і
$\Hfield=\Bfield$ дає систему двох рівнянь
\begin{equation}\label{eq:FresnelPBoundary}
	\cos i\,(E_m-E'_m) = \cos r\,E''_m, \qquad n_1(E_m+E'_m)=n_2E''_m.
\end{equation}
Виключаючи з другого рівняння $n_1=n_2\sin r/\sin i$ (закон Снелла) і підставляючи в нього $E''_m$ з першого рівняння, маємо
\begin{multline*}
	\sin r\cos r\,(E_m+E'_m) = \sin i\cos i\,(E_m-E'_m) \quad\Longrightarrow\\E_m\big[\sin2i-\sin2r\big] = E'_m\big[\sin2i+\sin2r\big]
\end{multline*}
(тут використано $\sin2\alpha=2\sin\alpha\cos\alpha$). Застосувавши до різниці й суми синусів подвійних кутів формули $\sin2i-\sin2r=2\cos(i+r)\sin(i-r)$
та $\sin2i+\sin2r=2\sin(i+r)\cos(i-r)$, знаходимо $E'_m/E_m=\tg(i-r)/\tg(i+r)$; підстановка цього результату в $E''_m=\dfrac{n_1}{n_2}(E_m+E'_m)=\dfrac{\sin
r}{\sin i}(E_m+E'_m)$ дає $E''_m$. Разом це --- \emphx{формули Френеля} для $p$-поляризації:
\begin{equation}\label{eq:FresnelPReflection}
	\tcbhighmath{
		E'_\parallel = \frac{\tg(i-r)}{\tg(i+r)}\, E_\parallel,
	}
\end{equation}
\begin{equation}\label{eq:FresnelPTransmission}
	\tcbhighmath{
		E''_\parallel = \frac{2\sin r\cos i}{\sin(i+r)\cos(i-r)}\, E_\parallel.
	}
\end{equation}

\begin{figure}[h!]\centering
	\localinput{p-polarization-geometry.tikz}
	\caption{Падаюча, відбита й заломлена $p$-поляризовані хвилі на межі розділу двох
		діелектриків\label{pic:PPolarization}}
\end{figure}

\begin{Attention}
	З формули~\eqref{eq:FresnelPReflection} випливає цікавий наслідок: якщо $i+r=90^\circ$, то $\tg(i+r)\to\infty$, і $E'_\parallel=0$ ---
	відбита хвиля \emphx{повністю відсутня} для $p$-поляризації! Кут падіння, за якого це відбувається, називають \emphx{кутом Брюстера}: за
	цієї умови відбите світло, незалежно від поляризації падаючого, виявляється \emphx{повністю $s$-поляризованим} --- явище, широко застосовуване
	в поляризаційній оптиці й фотографії (поляризаційні фільтри якраз і прибирають відблиски, поляризовані переважно за цим механізмом).
\end{Attention}

%% --------------------------------------------------------
\subsection*{Повне внутрішнє відбивання}
%% --------------------------------------------------------

Якщо світло переходить з оптично густішого середовища в оптично менш густе, $n_2<n_1$, то, за законом заломлення~\eqref{eq:SnellRefraction},
$\sin i/\sin r = n_2/n_1<1$, тобто $i<r$: кут заломлення завжди \emphx{більший} за кут падіння (рис.~\ref{pic:TotalReflection}).

\begin{figure}[h!]\centering
	\localinput{total-reflection-geometry.tikz}
	\caption{Заломлений промінь наближається до поверхні розділу при зростанні кута
		падіння\label{pic:TotalReflection}}
\end{figure}

Зі зростанням кута падіння $i$ заломлений промінь усе більше наближається до межі розділу. При деякому значенні $i=i_b$ кут заломлення
досягає $r=90^\circ$ --- на рисунку цьому відповідає промінь, що йде вздовж самої межі розділу. Кут $i_b$ називають \emphx{граничним
кутом}:
\begin{equation}\label{eq:CriticalAngle}
	\tcbhighmath{
		\sin i_b = \frac{n_2}{n_1} = n_{21}, \qquad\text{або}\qquad n\sin i_b=1,\ \ n=\frac{n_1}{n_2}>1.
	}
\end{equation}

При кутах падіння $i>i_b$ виникає \emphx{повне внутрішнє відбивання}: уся енергія падаючої хвилі відбивається назад у перше середовище, а
проходження заломленої хвилі в геометрично-оптичному сенсі стає неможливим.


%% --------------------------------------------------------
\subsection*{Що відбувається за граничним кутом: згасаюча хвиля}
%% --------------------------------------------------------

\enquote{Неможливість} існування заломленої хвилі при $i>i_b$ --- це висновок саме геометричної оптики. З'ясуємо, що відбувається із заломленою
хвилею формально, виходячи з \emphx{рівнянь Максвелла}. Заломлена хвиля $\Efield'' = \Efield''_m\, e^{i(\omega t-k''_xx-k''_yy)}$ (частота
$\omega''=\omega$ --- за~\eqref{eq:BoundaryConditionE}) повинна задовольняти дисперсійне співвідношення для другого середовища,
$k''^2={k''_x}^2+{k''_y}^2$, де $k''=n_2\omega/c$. З граничної умови $k_y=k''_y$~\eqref{eq:BoundaryConditionE} (а $k_y=k\sin i$, де
$k=n_1\omega/c$ --- хвильове число падаючої хвилі в першому середовищі) маємо
\begin{equation*}
	{k''_x}^2 = k''^2-k_y^2 = \left(\frac{n_2\omega}c\right)^{\!2}-\left(\frac{n_1\omega}c\right)^{\!2}\sin^2i = \frac{\omega^2}{c^2}\left(n_2^2-n_1^2\sin^2i\right).
\end{equation*}
Оскільки при $i>i_b$ виконується $n_1\sin i>n_2$ (тобто й $n_1^2\sin^2i>n_2^2$), вираз у дужках --- \emphx{від'ємний}, тому ${k''_x}^2<0$, тобто
$k''_x$ --- \emphx{уявне} число: $k''_x = iK$, де, повертаючись до позначення $n=n_1/n_2$ з формули~\eqref{eq:CriticalAngle},
\begin{equation*}
	K = \frac\omega c\sqrt{n_1^2\sin^2i-n_2^2} = \frac{n_2\omega}c\sqrt{n^2\sin^2i-1} > 0.
\end{equation*}
Розв'язок рівнянь Максвелла в другому середовищі при цьому набуває вигляду
\begin{equation}\label{eq:EvanescentWave}
	\tcbhighmath{
		\Efield'' = \Efield''_m\, e^{-Kx}\, e^{i(\omega t-k''_yy)}.
	}
\end{equation}
Отже, поле в другому середовищі \emphx{існує} навіть при куті падіння, більшому за граничний: воно \enquote{чіпляється} за межу розділу
й проникає в друге середовище на глибину порядку довжини хвилі світла $\lambda_0$, а амплітуда цього поля \emphx{експоненційно спадає}
вглиб середовища (уздовж $x$). Таке поле називають \emphx{згасаючою хвилею} (рис.~\ref{tikz:EvanescentWave}).

\begin{figure}[h!]\centering
\localinput{evanescent.tikz}
\caption{Згасаюча хвиля в другому середовищі при $i>i_b$ ($s$-поляризація: $\Efield=E_z\vect e_z$, спрямоване перпендикулярно до площини рисунка). Крива неперервно переходить із середовища $n_1$ (де поле --- суперпозиція падаючої й відбитої хвиль) у середовище $n_2$, де воно поширюється вздовж межі розділу (у напрямку $\vect{k}''$), а його амплітуда $|E_z|$ експоненційно спадає вглиб середовища (уздовж $x$) за законом $A(x)=A_0e^{-Kx}$; $x=1/K$ --- характерна глибина проникнення поля, порядку довжини хвилі $\lambda_0$.\label{tikz:EvanescentWave}}
\end{figure}

При $i\to i_b$ маємо $K\to0$, тобто глибина проникнення $1/K\to\infty$ --- саме тому граничний промінь на рис.~\ref{pic:TotalReflection}
($i=i_b$, $r=90^\circ$) виглядає як звичайний біжучий промінь уздовж межі розділу. Це узгоджується з уже отриманими формулами Френеля
для проходження, \eqref{eq:FresnelSTransmission} і~\eqref{eq:FresnelPTransmission}: підставляючи $r\to90^\circ$, маємо $E''_\perp\to
2E_\perp$ і $E''_\parallel\to2n\,E_\parallel$ --- обидві границі скінченні й ненульові, тобто поле на межі при $i=i_b$ цілком реальне.

Переконаємось прямим розрахунком, що нормальна до межі складова потоку Пойнтінга загасаючої хвилі дорівнює нулю в середньому. Для
$s$-поляризації ($\Efield''=E''_z\vect e_z$) магнітне поле знаходимо із~\eqref{eq:MaxwellPlaneWave}, $\vect k\times\Efield=\dfrac\omega
c\Bfield$, справедливого й для комплексного $\vect k''=(iK,k''_y,0)$:
\begin{equation*}
	\Bfield'' = \frac{c}{\omega}\,\vect k''\times\Efield'' = \frac{c}{\omega}\,E''_z\,(k''_y,\,-iK,\,0)
	\quad\Longrightarrow\quad
	B''_y = \frac{c}{\omega}k''_y E''_z, \qquad B''_x = -i\frac{c}{\omega}K E''_z.
\end{equation*}
Усереднений вектор Пойнтінга~\eqref{eq:PoyntingWaveDerivation} для двох полів однакової частоти дає $\langle\Pi_x\rangle\propto
-\mathrm{Re}(E''_z {H''_y}^{\!*})$, $\langle\Pi_y\rangle\propto\mathrm{Re}(E''_z{H''_x}^{\!*})$. Оскільки коефіцієнт при $E''_z$ у
$H''_y$ --- \emphx{уявний}, а в $H''_x$ --- \emphx{дійсний}:
\begin{equation}\label{eq:NormalFluxZero}
	\tcbhighmath{
		\langle\Pi_x\rangle = 0,
	} \qquad
	\langle\Pi_y\rangle \propto \frac{c}{\omega}k''_y\,|E''_m|^2\,e^{-2Kx} \neq 0.
\end{equation}

Звідки видно, що потік $\langle\Pi_y\rangle\ne0$ показує поширення енергії хвилі вздовж межі.


\begin{Attention}
Якщо на відстані порядку $\lambda_0$ від межі розділу розташувати третє середовище (наприклад, наблизити впритул ще один шматок скла до межі повного внутрішнього відбивання), \enquote{хвіст} загасаючої хвилі виявиться там ще достатньо відчутним, щоб розхитати електрони цього середовища й породити нову, вже \emphx{біжучу} хвилю --- частина світла \emphx{пройде} крізь проміжок, попри те, що кожна з двох меж окремо давала б повне відбивання. Це явище називають \emphx{порушеним повним внутрішнім відбиванням} --- оптичний аналог квантовомеханічного тунелювання, який легко спостерігати експериментально для сантиметрових радіохвиль (де $\lambda_0$ --- сантиметри, тому проміжок легко зробити порівнянним з довжиною хвилі) і значно важче --- для видимого світла (де $\lambda_0\sim10^{-5}$~см вимагає надзвичайно вузького проміжку).
\end{Attention}

%%% --------------------------------------------------------
%\subsection*{Що відбувається за граничним кутом: згасаюча хвиля}
%%% --------------------------------------------------------
%
%\enquote{Неможливість} існування заломленої хвилі при $i>i_b$ --- це висновок саме геометричної оптики. З'ясуємо, що відбувається із заломленою
%хвилею формально, виходячи з \emphx{рівнянь Максвелла}. Заломлена хвиля $\Efield'' = \Efield''_m\, e^{i(\omega t-k''_xx-k''_yy)}$ (частота
%$\omega''=\omega$ --- за~\eqref{eq:BoundaryConditionE}) повинна задовольняти дисперсійне співвідношення для другого середовища,
%$k''^2={k''_x}^2+{k''_y}^2$, де $k''=n_2\omega/c$. З граничної умови $k_y=k''_y$~\eqref{eq:BoundaryConditionE} (а $k_y=k\sin i$, де
%$k=n_1\omega/c$ --- хвильове число падаючої хвилі в першому середовищі) маємо
%\begin{equation*}
%	{k''_x}^2 = k''^2-k_y^2 = \left(\frac{n_2\omega}c\right)^{\!2}-\left(\frac{n_1\omega}c\right)^{\!2}\sin^2i = \frac{\omega^2}{c^2}\left(n_2^2-n_1^2\sin^2i\right).
%\end{equation*}
%Оскільки при $i>i_b$ виконується $n_1\sin i>n_2$ (тобто й $n_1^2\sin^2i>n_2^2$), вираз у дужках --- \emphx{від'ємний}, тому ${k''_x}^2<0$, тобто
%$k''_x$ --- \emphx{уявне} число: $k''_x = iK$, де, повертаючись до позначення $n=n_1/n_2$ з формули~\eqref{eq:CriticalAngle},
%\begin{equation*}
%	K = \frac\omega c\sqrt{n_1^2\sin^2i-n_2^2} = \frac{n_2\omega}c\sqrt{n^2\sin^2i-1} > 0.
%\end{equation*}
%Розв'язок рівнянь Максвелла в другому середовищі при цьому набуває вигляду
%\begin{equation}\label{eq:EvanescentWave}
%	\tcbhighmath{
%		\Efield'' = \Efield''_m\, e^{-Kx}\, e^{i(\omega t-k''_yy)}.
%	}
%\end{equation}
%Отже, поле в другому середовищі \emphx{існує} навіть при куті падіння, більшому за граничний: воно \enquote{чіпляється} за межу розділу
%й проникає в друге середовище на глибину порядку довжини хвилі світла $\lambda_0$, а амплітуда цього поля \emphx{експоненційно спадає}
%вглиб середовища (уздовж $x$). Таке поле називають \emphx{згасаючою хвилею} (рис.~\ref{tikz:EvanescentWave}).
%
%\begin{figure}[h!]\centering
%\localinput{evanescent.tikz}
%\caption{Згасаюча хвиля в другому середовищі при $i>i_b$ ($s$-поляризація: $\Efield=E_z\vect e_z$, спрямоване перпендикулярно до площини рисунка). Крива неперервно переходить із середовища $n_1$ (де поле --- суперпозиція падаючої й відбитої хвиль) у середовище $n_2$, де воно поширюється вздовж межі розділу (у напрямку $\vect{k}''$), а його амплітуда $|E_z|$ експоненційно спадає вглиб середовища (уздовж $x$) за законом $A(x)=A_0e^{-Kx}$; $x=1/K$ --- характерна глибина проникнення поля, порядку довжини хвилі $\lambda_0$.\label{tikz:EvanescentWave}}
%\end{figure}
%
%
%При $i\to i_b$ маємо $K\to0$, тобто глибина проникнення $1/K\to\infty$: саме тому граничний промінь на рис.~\ref{pic:TotalReflection}
%($i=i_b$, $r=90^\circ$) виглядає як звичайний біжучий промінь уздовж межі розділу.
%
%Це узгоджується й з уже отриманими формулами Френеля для проходження, \eqref{eq:FresnelSTransmission} і~\eqref{eq:FresnelPTransmission}.
%Підставляючи в них $r\to90^\circ$ (тобто $\sin r\to1$, $\sin(i_b+r)\to\cos i_b$, $\cos(i_b-r)\to\sin i_b$), маємо
%\begin{equation*}
%	E''_\perp \to 2E_\perp, \qquad E''_\parallel \to \frac{2}{\sin i_b}\,E_\parallel \overset{\eqref{eq:CriticalAngle}}{=} 2n\,E_\parallel.
%\end{equation*}
%Переконаємось прямим розрахунком, що нормальна до межі складова потоку Пойнтінга загасаючої хвилі справді дорівнює нулю в середньому.
%Візьмемо для означеності $s$-поляризацію, $\Efield''=E''_z\vect e_z$. Магнітне поле знаходимо з уже отриманого
%співвідношення~\eqref{eq:MaxwellPlaneWave}, $\vect k\times\Efield=\dfrac\omega c\Bfield$ --- воно справедливе для \emphx{будь-якого}
%хвильового вектора, зокрема й комплексного $\vect k''=(iK,\, k''_y,\, 0)$:
%\begin{equation*}
%	\Bfield'' = \frac{c}{\omega}\,\vect k''\times\Efield'' = \frac{c}{\omega}\,E''_z\,(k''_y,\,-iK,\,0),
%\end{equation*}
%звідки тангенційна складова $B''_y = \dfrac{c}{\omega}k''_y E''_z$ (дійсний коефіцієнт при $E''_z$), а нормальна ---
%$B''_x = -i\dfrac{c}{\omega}K E''_z$ (уявний коефіцієнт при $E''_z$).
%
%Для двох полів однакової частоти, записаних у комплексній формі, середнє за період добутку їхніх дійсних частин дорівнює
%$\langle\mathrm{Re}(A)\,\mathrm{Re}(B)\rangle = \tfrac12\mathrm{Re}(AB^*)$. Тому усереднений вектор Пойнтінга
%$\vect\Pi=\dfrac{c}{4\pi}[\Efield\times\Hfield]$ (див.~\eqref{eq:PoyntingWaveDerivation}) набуває вигляду
%$\langle\vect\Pi\rangle = \dfrac{c}{8\pi}\mathrm{Re}\!\left[\Efield''\times{\Hfield''}^{\!*}\right]$. Для нашої геометрії ($\Efield''\parallel\vect
%e_z$) маємо $\Efield''\times{\Hfield''}^{\!*} = \left(-E''_z {H''_y}^{\!*},\ E''_z {H''_x}^{\!*},\ 0\right)$, тобто
%\begin{equation}\label{eq:PoyntingEvanescentNormal}
%	\langle\Pi_x\rangle \propto -\mathrm{Re}\!\left(E''_z {H''_y}^{\!*}\right), \qquad \langle\Pi_y\rangle \propto \mathrm{Re}\!\left(E''_z {H''_x}^{\!*}\right).
%\end{equation}
%Оскільки $H''_y = -i\dfrac{c}{\omega}K E''_z$, маємо ${H''_y}^{\!*} = i\dfrac{c}{\omega}K {E''_z}^{\!*}$, і
%\begin{equation*}
%	E''_z {H''_y}^{\!*} = i\,\frac{c}{\omega}K\,|E''_z|^2
%\end{equation*}
%--- це число \emphx{чисто уявне} (бо $K,c,\omega,|E''_z|^2$ --- дійсні), тому його дійсна частина дорівнює нулю:
%\begin{equation}\label{eq:NormalFluxZero}
%	\tcbhighmath{
%		\langle\Pi_x\rangle = 0.
%	}
%\end{equation}
%Натомість $H''_x = \dfrac{c}{\omega}k''_y E''_z$ (дійсний коефіцієнт), тому $E''_z{H''_x}^{\!*} = \dfrac{c}{\omega}k''_y|E''_z|^2$ ---
%величина \emphx{дійсна й додатна}, і
%\begin{equation*}
%	\langle\Pi_y\rangle \propto \frac{c}{\omega}k''_y\,|E''_m|^2\,e^{-2Kx} \neq 0.
%\end{equation*}
%
%\begin{Attention}
%	Формула~\eqref{eq:NormalFluxZero} --- строгий доказ того, що загасаюча хвиля не переносить енергії \emphx{в середньому} крізь межу
%	розділу (саме тому відбивання й називають \emphx{повним}). Причина --- \emphx{зсув фази на $90°$} між $E''_z$ і $H''_y$: множник $i$ у
%	$H''_y=-i(c/\omega)K E''_z$ означає, що нормальна складова магнітного поля відстає (чи випереджає) електричне на чверть періоду, а
%	добуток двох величин, зсунутих по фазі на $90°$, у середньому за період дає нуль --- це той самий механізм, що й для реактивної
%	потужності в змінному струмі. Тангенційна ж складова $H''_x$ синфазна з $E''_z$ (коефіцієнт дійсний), тому $\langle\Pi_y\rangle\ne0$:
%	енергія реально тече вздовж межі розділу. Саме цим тангенційним потоком пояснюється \emphx{зсув Гуса--Генхен} (Goos--Hänchen shift):
%	відбитий промінь при повному внутрішньому відбиванні виходить трохи зміщеним уздовж межі відносно положення, яке передбачає проста
%	геометрична оптика, --- ніби він встигає \enquote{пробігти} деяку відстань уздовж межі в товщі загасаючої хвилі, перш ніж повернутися
%	в перше середовище.
%\end{Attention}
%
%Якщо ж на відстані порядку $\lambda_0$ від межі розділу розташувати третє середовище (наприклад, наблизити впритул ще один шматок скла
%до межі повного внутрішнього відбивання), \enquote{хвіст} загасаючої хвилі виявиться там ще достатньо відчутним, щоб розхитати
%електрони цього середовища й породити нову, вже \emphx{біжучу} хвилю --- частина світла \emphx{пройде} крізь проміжок, попри те, що
%кожна з двох меж окремо давала б повне відбивання. Це явище називають \emphx{порушеним повним внутрішнім відбиванням} --- оптичний
%аналог квантовомеханічного тунелювання, який легко спостерігати експериментально для сантиметрових радіохвиль (де $\lambda_0$ ---
%сантиметри, тому проміжок легко зробити порівнянним з довжиною хвилі) і значно важче --- для видимого світла (де
%$\lambda_0\sim10^{-5}$~см вимагає надзвичайно вузького проміжку).

%% --------------------------------------------------------
\section*{Підсумки}
%% --------------------------------------------------------

%\subsection*{Хвильове рівняння}
%\begin{itemize}
%	\item Загальний вигляд хвилі: $u=f(x-vt)+g(x+vt)$; хвильове рівняння: $\nabla^2 u = \dfrac1{v^2}\pparttime u$.
%	\item Плоска монохроматична хвиля: $u=u_m\cos(\omega t-\vect k\cdot\vect r)$; $\lambda=2\pi/k$, $v=\omega/k$.
%\end{itemize}
%
%\subsection*{Електромагнітні хвилі}
%\begin{itemize}
%	\item З рівнянь Максвелла (без джерел): $\nabla^2\Efield = \dfrac1{v^2}\pparttime\Efield$, $\nabla^2\Hfield=\dfrac1{v^2}\pparttime\Hfield$,
%	      $v=c/\sqrt{\epsilon\mu}$; показник заломлення $n=\sqrt{\epsilon\mu}$.
%	\item Поперечність: $\Efield\perp\vect k$, $\Hfield\perp\vect k$, $\Efield\perp\Hfield$; амплітуди: $nE_m=B_m$, $\sqrt\epsilon E_m=\sqrt\mu
%	      H_m$.
%\end{itemize}
%
%\subsection*{Енергія, тиск, імпульс}
%\begin{itemize}
%	\item $\vect\Pi = w\vect v$; тиск хвилі: $p=w$ (поглинання), $p=w(1+\rho)$ (часткове відбивання).
%	\item Густина імпульсу поля: $\vect g = \vect\Pi/c^2$.
%\end{itemize}
%
%\subsection*{Відбивання й заломлення}
%\begin{itemize}
%	\item Закони відбивання й заломлення: $i=i'$, $n_1\sin i=n_2\sin r$.
%	\item Формули Френеля: $E'_\perp = -\dfrac{\sin(i-r)}{\sin(i+r)}E_\perp$, $E'_\parallel = \dfrac{\tg(i-r)}{\tg(i+r)}E_\parallel$.
%	\item Повне внутрішнє відбивання: $\sin i_b = n_2/n_1$; загасаюча хвиля $\propto e^{-Kx}$ при $i>i_b$.
%\end{itemize}

Розділом про електромагнітні хвилі замикається логічне коло, розпочате законом Кулона: рівняння Максвелла, виведені з дослідів про сили між
нерухомими зарядами й провідниками зі струмом, самі собою породжують хвильовий рух, що поширюється зі швидкістю світла, --- і саме світло, а
з ним і вся геометрична й хвильова оптика, виявляються лише \emphx{частинним випадком} електромагнітних явищ. У цьому --- найглибший
синтетичний висновок усього курсу: електрика, магнетизм і оптика --- не три окремі науки, а три прояви \emphx{одного} електромагнітного поля.


\begin{Questions}
\begin{quizlist}
	\item Що таке хвиля в загальному сенсі? Наведіть кілька прикладів фізичних процесів, що описуються хвилями, і поясніть, що означає величина $u$ в кожному з них.
	\item Виведіть хвильове рівняння $\partial^2u/\partial x^2 = (1/v^2)\,\partial^2u/\partial t^2$ для збурення $u=f(x-vt)+g(x+vt)$, що поширюється без зміни форми. Як узагальнюється це рівняння на тривимірний випадок?
	\item Чому рівняння $\nabla^2u = (1/v^2)\,\partial^2u/\partial t^2$ вважають \enquote{паролем} хвильового процесу в різних розділах фізики?
	\item Що таке плоска монохроматична хвиля? Запишіть її дійсну та комплексну форми. У чому зручність комплексного запису?
	\item Дайте означення довжини хвилі, хвильового вектора, хвильового числа, колової частоти, періоду та фазової швидкості хвилі. Як вони пов'язані між собою?
	\item Виведіть хвильові рівняння для $\Efield$ і $\Hfield$ з рівнянь Максвелла в однорідному середовищі без вільних зарядів і струмів. Чому дорівнює швидкість поширення електромагнітних хвиль і як з неї отримати показник заломлення $n$?
	\item Чому збіг швидкості електромагнітних хвиль зі швидкістю світла, отриманий суто з рівнянь електрики й магнетизму, вважають одним із найважливіших результатів класичної фізики?
	\item Запишіть рівняння Максвелла для плоскої гармонічної хвилі в комплексній формі. Яке правило заміни диференціальних операцій алгебраїчними при цьому використовується?
	\item Доведіть поперечність електромагнітних хвиль ($\Efield\perp\vect k$, $\Hfield\perp\vect k$). Яке співвідношення пов'язує амплітуди $E_m$ і $H_m$, і що з нього випливає для незалежності полів $\Efield$ і $\Hfield$?
	\item Покажіть, що електрична й магнітна частини густини енергії хвилі рівні між собою. Виведіть формулу $\vect\Pi = w\vect v$ для вектора Пойнтінга хвилі та поясніть її фізичний зміст.
	\item Що таке інтенсивність випромінювання і чому вона пропорційна $E_m^2$?
	\item Виведіть формулу тиску електромагнітної хвилі $p=w$ при повному поглинанні та $p=w(1+\rho)$ при частковому відбиванні. Чому при повному відбиванні тиск удвічі більший, ніж при повному поглинанні?
	\item Опишіть дослід Лебедєва: яку ідею він перевіряв і як влаштування крутильних терезів дозволило виміряти різницю тисків для поглинального й відбивального крилець?
	\item Яку роль відіграє тиск світла в астрофізиці (рівновага зір, хвости комет) і в космічній техніці (сонячні та лазерні вітрила)? Наведіть приклади місій, що використовують цей принцип.
	\item Виведіть формулу для густини імпульсу поля $\vect g = \vect\Pi/c^2$ двома способами: з релятивістського співвідношення енергія--імпульс для безмасової частинки та прямим розрахунком сили Ампера.
	\item У чому полягає парадокс Фейнмана з циліндричним конденсатором в аксіальному магнітному полі? Виведіть момент імпульсу поля конденсатора $L=qBR^2/(2c)$ та кутову швидкість його обертання після вимкнення поля.
	\item Як парадокс Фейнмана демонструє, що електромагнітне поле --- повноцінний фізичний об'єкт, здатний нести імпульс і момент імпульсу?
	\item Що таке $s$- та $p$-поляризація хвилі відносно площини падіння? Чому довільну поляризацію можна розкласти на суперпозицію цих двох випадків?
	\item Виведіть закони відбивання ($i=i'$) та заломлення ($n_1\sin i = n_2\sin r$) з граничної умови неперервності тангенціальної складової $\Efield$ для $s$-поляризованої хвилі. Чому рівність частот $\omega=\omega'=\omega''$ означає незмінність кольору світла при відбиванні й заломленні?
	\item Чому закони геометричної оптики (відбивання й заломлення) вважають прямим наслідком граничних умов для електромагнітного поля, а не окремим постулатом?
	\item Запишіть формули Френеля для $s$-поляризації. Що вони визначають?
	\item Запишіть формули Френеля для $p$-поляризації. Що таке кут Брюстера і чому за цієї умови відбите світло виявляється повністю $s$-поляризованим?
	\item Що таке граничний кут повного внутрішнього відбивання? Виведіть умову $\sin i_b = n_2/n_1$ і поясніть, чому вона можлива лише при переході з оптично густішого середовища в менш густе.
	\item Що відбувається з полем у другому середовищі при куті падіння, більшому за граничний? Виведіть вигляд загасаючої (evanescent) хвилі та поясніть, чому вона не переносить енергії в середньому вглиб середовища.
	\item Що таке порушене повне внутрішнє відбивання (frustrated total internal reflection)? Чому це явище вважають оптичним аналогом квантовомеханічного тунелювання, і чому його легше спостерігати для радіохвиль, ніж для видимого світла?
\end{quizlist}
\end{Questions}